% NSF Data Management and Sharing Plan -- CONTENT DRAFT for the Research.gov webform
%
% Since April 27, 2026 (PAPPG 24-1, Supplement 2), the plan is no longer uploaded as a PDF:
% it is entered in a Research.gov webform generated for the proposal's directorate (MPS).
% NSF 23-603's instruction to upload a 2-page "Data Management Plan" predates this change.
% For each research-product category (at most four), the form asks for: general information
% (category, title, description); access policies and limitations; data standards and
% metadata; provenance; public archiving; timeline for public accessibility; data
% availability (retention period, plus confirmation that the repository's retention policy
% will be followed); and accountability (responsible PI). The fields below follow that
% order, and the first words of each field are the closest dropdown choice where the form
% has one. Opting out requires a justification that the project produces no data or research
% products, which is not the case here.
%
% The PI needs no computing (2026-10-20), so there is no software category. NSF's guidance
% warns against using GitHub as an archive, so arXiv is the archive for papers and GitHub
% only mirrors their LaTeX sources.

\documentclass[11pt]{article}

\usepackage[margin=1in]{geometry}
\usepackage{amsmath, amssymb}
\usepackage{microtype}
\usepackage{enumitem}
\usepackage[compact]{titlesec}

\titleformat{\section}{\large\bfseries}{\thesection.}{0.5em}{}
\titlespacing*{\section}{0pt}{1.2ex plus .4ex minus .2ex}{0.5ex plus .2ex}
\setlist{nosep, leftmargin=*}
\setlist[description]{font=\normalfont\itshape, leftmargin=0pt}

% Block paragraphs, as in the Project Description.
\setlength{\parindent}{0pt}
\setlength{\parskip}{0.3\baselineskip plus 2pt minus 1pt}

\pagestyle{empty}

\begin{document}

\begin{center}
  {\Large\bfseries Data Management and Sharing Plan}\\[5pt]
  \emph{Content for the Research.gov webform, which has replaced the PDF upload.}
\end{center}

This project is theoretical research in pure mathematics. It involves no human subjects,
sensitive data, or physical samples or collections, and it requires no computation. Its
research products are research articles and expository materials from the Broader Impacts
activities. All of them will be shared openly, free of charge and without restriction.

\section*{Category 1: Research articles and preprints}

\begin{description}
  \item[Description.] Research papers presenting the results of the project (theorems,
  proofs, and examples), and an expository survey if one is written.
  \item[Access policies and limitations.] Not applicable: there are no restrictions on
  access.
  \item[Data standards and metadata.] PDF, with \LaTeX\ source deposited on arXiv; arXiv
  metadata (title, authors, abstract, and Mathematics Subject Classification codes); DOIs
  for published versions.
  \item[Provenance.] New (user-defined): original research by the PI.
  \item[Public archiving.] arXiv for preprints, and the NSF Public Access Repository
  (NSF-PAR) for the accepted manuscript or published version of each journal article. The
  \LaTeX\ sources will also be posted in a public GitHub repository.
  \item[Timeline for public accessibility.] At time of publication. Each preprint will be
  posted to arXiv no later than its submission to a journal, and each accepted manuscript
  will be deposited in NSF-PAR when it is published, in accordance with NSF's public access
  policy.
  \item[Data availability.] Indefinitely. arXiv is a permanent archive, and NSF-PAR provides
  long-term public access; the PI will follow both repositories' retention policies.
  \item[Accountability.] Edward Varvak (PI).
\end{description}

\section*{Category 2: Seminar and workshop materials}

\begin{description}
  \item[Description.] Notes and recorded talks from the learning seminar on Lagrangian
  fibrations and Hitchin systems, and lecture materials from the early-career workshop,
  both described in the Broader Impacts section of the Project Description.
  \item[Access policies and limitations.] Not applicable. Talks will be recorded and posted
  only with the speaker's consent.
  \item[Data standards and metadata.] PDF (notes typeset in \LaTeX) and MP4 video, each
  labeled with title, speaker, and date.
  \item[Provenance.] New (user-defined): materials created by the PI and seminar speakers.
  \item[Public archiving.] Zenodo, with links from the seminar's web page.
  \item[Timeline for public accessibility.] User-defined: within one month of each seminar
  meeting or workshop.
  \item[Data availability.] At least 10 years, following Zenodo's retention policy.
  \item[Accountability.] Edward Varvak (PI).
\end{description}

\end{document}
